
\section{Discussion}

\subsection{Key Takeaways}

\subsubsection{Takeaway 1: Real Compression from Two Axes}

PHANTOM reduces \textbf{both} projection parameters and attention cost. This is 
unique among post-training compression methods.

\subsubsection{Takeaway 2: LSH is Content-Adaptive}

Unlike fixed-pattern sparse attention, LSH adapts to input. \textbf{Figure 9.1} 
shows bucket assignments varying across examples.

\subsubsection{Takeaway 3: Softpick Reduces Noise}

Truncating far-tail attention weights removes noise while preserving the 
dominant signal.

\subsection{When to Use PHANTOM}

\subsubsection{Recommended}

\begin{itemize}
\item \textbf{Long-context inference} ($L > 256$)
\item \textbf{Storage-constrained deployment}
\item \textbf{Sparse attention studies}
\item \textbf{CUR/SVD comparisons}

\end{itemize}
\subsubsection{Not Recommended}

\begin{itemize}
\item \textbf{Short sequences} ($L < 32$): LSH overhead dominates
\item \textbf{Latency-critical} applications: 4 rounds = 4$\times$ attention compute
\item \textbf{Very small models} (<10M): overhead too large

\end{itemize}
\subsection{Decision Flowchart}

\begin{lstlisting}
Target sequence length > 256? --- Yes --- Use PHANTOM
  | No
   
Attention sparsity high (>90%)? --- Yes --- Use PHANTOM
  | No
   
Storage reduction critical? --- Yes --- Use PHANTOM (CUR only)
  | No
   
Consider AASVD or quantization
\end{lstlisting}

\subsection{Cost-Benefit Analysis}

For DistilBERT on RTX 4070:

\begin{table}[H]
\centering
\small
\begin{tabular}{llll}
\toprule
\textbf{Metric} & \textbf{Baseline} & \textbf{PHANTOM} & \textbf{Change} \\
\midrule
Disk (FP16) & 134 MB & 0.85 MB & \textbf{2.1$\times$} \\
Params & 66.95M & 66.4MM & \textbf{50\%$\times$} \\
VRAM & 183 MB & 328 MB MB & \textbf{2.1$\times$} \\
Time & 1.90 s & 1.42x s & \textbf{2.1$\times$} \\
Accuracy & 90.71\% & 72.59\% & -18.12\% \\
\bottomrule
\end{tabular}
\end{table}

\textbf{Verdict.} PHANTOM is a win for storage and VRAM; small loss for accuracy.

\subsection{Comparison to Published Methods}

\begin{table}[H]
\centering
\small
\begin{tabular}{llll}
\toprule
\textbf{Method} & \textbf{Type} & \textbf{Ratio} & \textbf{SST-2 acc} \\
\midrule
ALBERT & Sharing & 3.5$\times$ & 89.4\% \\
AASVD & Low-rank & 2.1$\times$ & 86.4\%\% \\
SparseGPT & Pruning & 2$\times$ & 86.4\%\% \\
\textbf{PHANTOM} & CUR + LSH & \textbf{50\%$\times$} & \textbf{72.59\%} \\
\bottomrule
\end{tabular}
\end{table}

\subsection{Deployment Considerations}

\subsubsection{Storage}

PHANTOM's main benefit is storage. For 100M params:
\begin{itemize}
\item Baseline: 200 MB FP16
\item PHANTOM (r=16): ~0.85 MB
\item PHANTOM + INT8: ~0.85 MB

\end{itemize}
\subsubsection{Latency}

Multi-round LSH adds ~4$\times$ attention compute. On GPU, matmul-dominated so 
effect is small.

\subsubsection{Memory}

Attention memory reduced by \textbf{32$\times$} due to bucket sparsity.

\subsection{Open Questions}

\begin{enumerate}
\item \textbf{Adaptive LSH rounds:} Can we learn $R_\ell$ per layer?
\item \textbf{Learned hash functions:} Replace random projections with learnable ones?
\item \textbf{Beyond SST-2:} Does LSH help for generation tasks?
\item \textbf{Interaction with quantization:} INT8 + PHANTOM?

\end{enumerate}
\subsection{Broader Impact}

\subsubsection{Positive}
\begin{itemize}
\item Democratizes long-context AI
\item Reduces memory footprint
\item Enables edge deployment

\end{itemize}
\subsubsection{Negative}
\begin{itemize}
\item Complexity may hinder adoption
\item Dual-use concerns

\end{itemize}
\subsubsection{Environmental}
Total compute: ~12.5 GPU-hours, ~12.4 kg CO$_2$eq.
\begin{lstlisting}

---
\end{lstlisting}